% concatenated from mcc.tex
\documentclass[12pt]{amsart}

%packages
\usepackage[utf8]{inputenc}
\usepackage{amsmath}
\usepackage{amsthm}
\usepackage{amsfonts}
\usepackage{amssymb}
\usepackage{tabu}
\usepackage{latexsym} 
\usepackage[margin=1in]{geometry}
\usepackage{enumerate} 
\usepackage[colorlinks=true]{hyperref}
\usepackage{mathrsfs}
\usepackage{tikz}
\usepackage[parfill]{parskip}
\setlength{\parindent}{0pt}
\usetikzlibrary{matrix,arrows,decorations.pathmorphing}
\usepackage{tikz-cd}
\usepackage[all]{xy}
%\usepackage{titlesec} % Allows customization of titles
\usepackage{setspace}
\usepackage{tabularx}
\usepackage{xcolor}
\usepackage{framed}


%\usepackage{mathtools}
%Ext, Hom, End
\newcommand{\Ext}{\mathrm{Ext}}
\newcommand{\sExt}{\underline{\mathrm{Ext}}}
\newcommand{\Hom}{\mathrm{Hom}}
\newcommand{\End}{\mathrm{End}}
\newcommand{\sHom}{\underline{\mathrm{Hom}}}
\newcommand{\sEnd}{\underline{\mathrm{End}}}
\newcommand{\Tor}{\mathrm{Tor}}

%Ideals
\newcommand{\ann}{\mathrm{ann}}
\newcommand{\sann}{\underline{\mathrm{ann}}}
\newcommand{\tsann}{\mathrm{tot}\sann}
\newcommand{\arann}{\mathrm{ARann}}
\newcommand{\jac}{\mathrm{jac}}
\newcommand\ca{\mathrm{ca}}
\newcommand\co{\mathrm{co}}
\newcommand{\m}{\mathfrak{m}}
\newcommand{\E}{\mathsf{E}}
\newcommand{\TTT}{\mathbb{T}}
%Categories
\newcommand\TT{\mathcal{T}}
\newcommand{\module}{\mathrm{mod}}
\newcommand{\smodule}{\underline{\mathrm{mod}}}
\newcommand{\MCM}{\mathrm{MCM}}
\newcommand{\ZCM}{\mathrm{ZCM}}
\newcommand{\sMCM}{\underline{\mathrm{MCM}}}
\newcommand{\CM}{\mathrm{CM}}
\newcommand{\sCM}{\underline{\mathrm{CM}}}
\newcommand{\add}{\mathrm{add}}
\newcommand{\gen}{\mathrm{gen}}
\newcommand{\cogen}{\mathrm{cogen}}
\newcommand{\proj}{\mathrm{proj}}
\newcommand{\RI}{\mathcal{RI}}
\newcommand{\CMI}{\CM{\text{-}}\mathcal{I}}
\newcommand{\reflexive}{\mathrm{ref}}
\newcommand{\Ul}{\mathrm{Ul}}



\DeclareMathOperator{\kac}{\mathrm{KAC}}

%Numerical Invariants
\newcommand{\depth}{\mathrm{depth }}
\newcommand{\gldim}{\mathrm{gldim }}
\newcommand{\pd}{\mathrm{pd }}
\newcommand{\CMdd}{\mathrm{ZCMdomdim }}
\newcommand{\id}{\mathrm{id }}
\newcommand{\CMid}{\mathrm{CMid }}
\newcommand{\rd}{\mathrm{rd }}
\newcommand{\CMrd}{\mathrm{CMrigd }}
\newcommand{\radius}{\mathrm{radius}}

%Number Sets
\newcommand{\CC}{\mathcal{C}}
\newcommand{\C}{\mathbb{C}}
\newcommand{\ZZ}{\mathbb{Z}}
\newcommand{\QQ}{\mathbb{Q}}
\newcommand{\DD}{\mathcal{D}}
\newcommand{\PP}{\mathbb{P}}
\newcommand{\XX}{\mathcal{X}}
\newcommand{\PPP}{\mathcal{P}}
\newcommand{\EE}{\mathcal{E}}
%Extra
\newcommand{\spec}{\mathsf{Spec}}
\newcommand{\sing}{\mathrm{Sing}}
\newcommand{\st}{\mathrm{st}}
\newcommand{\coker}{\mathrm{coker}}
\newcommand{\Lam}{\Lambda}
\newcommand{\bosluk}{\\ \phantom{x} \\}
\newcommand{\imof}{\mathrm{im}}
\newcommand{\FF}{\mathbb{F}}
\newcommand{\KK}{\mathbb{K}}
\newcommand{\bfa}{\mathbf{a}}
\newcommand{\bfb}{\mathbf{b}}
\newcommand{\om}{\Omega}
\newcommand{\si}{\Sigma}
%Theorems
\newtheorem{theorem}{Theorem}[section]
\newtheorem*{theorem*}{Theorem}
\newtheorem{lemma}[theorem]{Lemma}
\newtheorem*{lemma*}{Lemma}
\newtheorem{proposition}[theorem]{Proposition}
\newtheorem*{proposition*}{Proposition}
\newtheorem{corollary}[theorem]{Corollary}
\newtheorem*{corollary*}{Corollary}
\newtheorem{question}[theorem]{Question}
\theoremstyle{definition}
\newtheorem{definition}[theorem]{Definition}
\newtheorem*{definition*}{Definition}
\newtheorem{example}[theorem]{Example}
\newtheorem{remark}[theorem]{Remark}
\newtheorem{exercise}[theorem]{Exercise}
\newtheorem{conjecture}[theorem]{Conjecture}
\newtheorem{chunk}[theorem]{}
\newtheorem{setting}[theorem]{Setting}




\newtheorem{introtheorem}{Theorem}
\renewcommand{\theintrotheorem}{\Alph{introtheorem}}

\newcommand{\Ozgur}{\"{O}zg\"{u}r }
%
%Section Formatting
%\titleformat{\chapter}[block]{\Huge\scshape\bf\centering}{\thechapter.}%{1em}{} % Change the look of the section titles
%\titleformat{\section}[block]{\large\scshape\bf\centering}{\thesection.}%{1em}{} % Change the look of the section titles
%\titleformat{\subsection}[block]{\scshape\bf\centering}{\thesubsection.}{1em}{} % Change the look of the section titles

\DeclareMathOperator{\tr}{\mathrm{tr}}
\newcommand\oznewtext[1]{{\color{olive} #1}} %new text Ozgur
\newcommand\ozcomment[1]{{\color{magenta} \sf {\"O}zg{\"u}r: [#1]}}
\newcommand\ozoldtext[1]{{\color{red} #1}}


%\input{~/preambles/preamble.tex}

\title{When is the canonical conductor minimal?}
\author[\"{O}.~Esentepe]{\"{O}zg\"{u}r Esentepe}
\address{Institut für Mathematik und Wissenschaftliches Rechnen, Universität Graz,
Heinrichstraße 36, 8010 Graz, Austria}
\email{ozgur.esentepe@uni-graz.at}
\urladdr{https://www.sntp.ca}

\subjclass[2020]{13C05, 13C14, 13G05, 13H10}
\keywords{almost Gorenstein rings, nearly Gorenstein rings, far-flung Gorenstein rings, trace ideals, conductor ideals}

 
\begin{document}


\begin{abstract}
   For a one dimensional analytically unramified Cohen-Macaulay local ring $R$, the blowup algebra of the canonical ideal is a module finite birational extension. The conductor of this extension always contains the conductor of $R$. We study the case when there is equality. This is the case where $R$ is far from being almost Gorenstein. We study this property within the landscape of numerical semigroup rings and local Arf rings.
\end{abstract}

\maketitle
\thispagestyle{empty}

\section{Introduction}
Gorenstein rings are ubiquitous as Hyman Bass pointed out over half a century ago \cite{Bass}. They possess beautiful symmetry properties and they are a special family inside the broader family of Cohen-Macaulay rings. In the last few decades, there has been a growing interest in understanding Cohen-Macaulay local rings which are close to being Gorenstein. This has started with the introduction of one dimensional \textit{almost Gorenstein} rings by Barucci-Fr{\"o}berg \cite{Barucci-Froberg} and continued with the work of Goto-Matsuoka-Phuong \cite{Goto-Matsuoka-Phuong} and Goto-Takahashi-Taniguchi \cite{Goto-Takahashi-Taniguchi} in higher dimensions. Later in \cite{Herzog-Hibi-Stamate}, Herzog-Hibi-Stamate introduced \textit{nearly Gorenstein} rings and this is still a very active research are in commutative algebra. On the other extreme, Herzog-Kumashiro-Stamate \cite{Herzog-Kumashiro-Stamate} considered \textit{far-flung-Gorenstein} rings which are one dimensional rings that are not close to being Gorenstein.

The main actors in this type of study are trace ideals. Nearly Gorenstein rings and far-flung Gorenstein rings are defined by imposing conditions on the trace of the canonical module. Dao-Maitra-Sridhar showed in \cite[Proposition 7.1]{Dao-Maitra-Sridhar} that in dimension one, the almost Gorenstein property can also be seen by a condition on the trace of the blowup algebra of the canonical module. More precisely, let $R$ be a one dimensional generically Gorenstein local ring with canonical ideal $\omega$ and maximal ideal $\m$. Denote by $\tr(\omega)$ the trace of the canonical module (see Definition \ref{trace-and-conductor}) and by $b(\omega)$ the conductor of the blowup algebra $B(\omega)$ of the canonical module (see Remark \ref{trace-conductor-remark} and Remark \ref{blow-up-remarks}). Then, under mild conditions there is a chain of ideals 
\begin{align*}
\co(R) \subseteq b(\omega) \subseteq \tr(\omega) \subseteq \m
\end{align*}
where $\co(R)$ is the conductor of $R$, provided that $R$ is not Gorenstein (in which case we have $b(\omega) = \tr(\omega) = R$ ). The three extreme cases have been studied:
\begin{enumerate}
    \item $R$ is a nearly Gorenstein ring if and only if the equality $\tr(\omega) = \m$ holds.
    \item $R$ is an almost Gorenstein ring if and only if the equality $b(\omega) = \m$ holds.
    \item $R$ is a far-flung Gorenstein ring if and only if the equality $\tr(\omega) = \co(R)$ holds.
\end{enumerate}
In this paper, we are treating the missing case: the case where the equality $b(\omega) = \co(R)$ holds. We say that such a ring has \textit{minimal canonical conductor}. By definition, the class of such rings contain the class of far-flung Gorenstein rings. Under some conditions, they coincide. For example, if $R$ has multiplicity 3, then $R$ has minimal canonical conductor if and only if it is far-flung Gorenstein. But in general the class of rings with minimal canonical conductor is much larger.

In Section 2, we recall some preliminaries and equivalent conditions which we use to define a ring with minimal canonical conductor. In Section 3, we compare this property with being far-flung Gorenstein and also study this property for blowups at the maximal ideal.

Section 4 compares having minimal canonical conductor with having reflexive and Gorenstein birational extensions. We then apply this to understand local Arf domains which have minimal canonical conductor.

Finally, in Section 5 we classify numerical semigroup rings which have minimal canonical conductor and we give an example which has minimal canonical conductor but is nearly Gorenstein.

\subsection*{Acknowledgements} I would like to thank Alfred Geroldinger, Alessio Moscariello and Shinya Kumashiro for valuable conversations and helpful comments. I also want to thank Souvik Dey for his comments after the first preprint version of this paper.


\section{Preliminaries}
We start this section with some conventions and general definitions. 
\begin{definition}\label{trace-and-conductor}
    Let $R$ be a commutative Noetherian local ring .
     \begin{enumerate}
        \item Let $M$ be a finitely generated $R$-module. Then, the \textit{trace ideal} $\tr(M)$ of $M$ is the ideal of $R$ generated by the elements of the form $f(m)$ where $f \in M^*$ and $m \in M$. 
        \item Let $Q(R)$ be the total quotient ring of $R$. By $\overline{R}$, we denote the normalisation of $R$ inside $Q(R)$. That, is $\overline{R}$ is the integral closure of $R$ inside its total quotient rings. Then, the \textit{conductor} ideal is defined as 
        \begin{align*}
            \co(R) = \{r \in \overline{R} \colon r \overline{R} \subseteq R \}.
        \end{align*}
        More generally, for any birational extension $R \subseteq S$, one can define 
        \begin{align*}
            \co(R,S) = \{r \in S \colon r S \subseteq R \}.
        \end{align*}
        Here, by a \textit{birational extension} we mean a ring extension $R \subseteq S \subseteq Q(R)$.
     \end{enumerate}
\end{definition}
\begin{remark}\label{trace-conductor-remark}
    Conductor ideals can be seen as trace ideals. This follows from \cite[Proposition 2.4]{Kobayashi-Takahashi-Trace} and has been carefully treated in \cite[Appendix]{Herzog-Hibi-Stamate-semigroup}. More precisely, one has $\co(R,S) = \tr(S)$. 
\end{remark}
Recall that a local ring is called \textit{analytically unramified} if its completion is reduced. For an analytically unramified local ring, the normalisation $\overline{R}$ is finitely generated as an $R$-module. The converse is also true when $R$ is a one dimensional Cohen-Macaulay local ring. In this case, $R$ is reduced. 
\begin{remark}
    Let $R$ be a one dimensional Cohen-Macaulay local ring which is analytically unramified. Assume that for a finitely generated $R$-module M, the trace ideal $\tr(M)$ contains a nonzerodivisor and has a principal reduction. Then, there exists an inclusion $\co(R) \subseteq \tr(M)$ \cite[Corollary 3.6]{Dao-Maitra-Sridhar}. When $R$ has infinite residue field, the principal reduction condition is satisfied automatically.
\end{remark}
\iffalse
The following lemma was established by Faber \cite[Theorem 3.5]{Faber} in the reduced case and more generally by Dao-Maitra-Sridhar \cite[Theorem 2.9]{Dao-Maitra-Sridhar}.
\begin{lemma}\label{faber-lemma}
    Let $R$ be a commutative Noetherian ring, $S$ a finite birational extension and $M$ a finitely generated $R$-module. Consider the two statements:
    \begin{enumerate}
        \item $M$ is a module over $S$.
        \item There is an inclusion $\tr(M) \subseteq \co(R,S)$. 
    \end{enumerate}
    Then, (1) implies (2). When $M$ is reflexive, (2) implies (1), as well.
\end{lemma}
\oznewtext{Now, assume that $R$ is a one dimensional analytically unramified Cohen-Macaulay local ring. This gives us that $R$ is reduced and that $\overline{R}$ is finitely generated as an $R$-module. Hence, Lemma \ref{faber-lemma} tells us that if $M$ is a reflexive $R$-module with $\tr(M) \subseteq \co(R)$, then $M$ is a reflexive $\overline{R}$-module. Note that a module $M$ over $\overline{R}$ is reflexive if and only if it is projective as an $\overline{R}$-module as $\overline{R}$ is a one dimensional normal (and hence regular) ring. In particular, if $\overline{R}$ is a local ring, then $\tr(M)$ implies $M \cong \overline{R}^{\oplus r}$ for some $r$ as an $R$-module.}
This gives us the following lemma about reflexive modules with minimal traces.
\begin{lemma}\label{main-lemma}
    Let $R$ be a one dimensional analytically unramified \oznewtext{local ring }and $M$ a reflexive $R$-module. If there exists an equality $\tr(M) = \co(R)$, then we have an isomorphism 
    \begin{align*}
    M \oznewtext{\oplus K} \cong \overline{R}^{\oplus r}
    \end{align*}
    \oznewtext{for some positive integer $r$ and a finitely generated $R$-module $K$.}
\end{lemma}
\oznewtext{
\begin{remark}
    \begin{enumerate}
        \item Lemma \ref{main-lemma} concerns modules with minimal trace. The opposite case of \textit{full-trace modules} have been studies by {\c C}elikba{\c s}-{\c C}elikba{\c s}-Herzog-Kumashiro in \cite{Celikbas-Celikbas-Herzog-Kumashiro}.
        \item The conductor ideal $\co(R)$ is a stable trace ideal and a general version of Lemma \ref{main-lemma} with birational extensions with stable conductors has been given by Dey in \cite[Corollary 4.5]{Dey}.
    \end{enumerate}
\end{remark}
}
\fi
\begin{definition}
    Let $R$ be a commutative Noetherian ring and $I$ a regular ideal. Then, there is a filtration of endomorphism algebras 
    \begin{align*}
    R \subseteq I \colon I \subseteq I^2 \colon I^2 \subseteq \cdots I^n \colon I^n \subseteq \cdots \subseteq Q(A).
    \end{align*}
    We put 
    \begin{align*}
    R^I = \bigcup_{n \geq 0} I^n \colon I^n
    \end{align*}        
    and denote by $b(I)$ the conductor ideal $\co(R,R^I)$.
\end{definition}
\begin{remark}\label{blow-up-remarks}
    Assume that $R$ is a commutative Noetherian ring of dimension one and $I$ a regular ideal.
    \begin{enumerate}
        \item The ring $R^I$ coincides with the blowup algebra $B(I)$ at $I$. This was proved by Lipman in his influencial work \cite{Lipman} on Arf rings.
        \item We say that $I$ is a \textit{stable} ideal if $R^I = I \colon I$.
        \item The ideal $b(I)$ equals the trace ideal $\tr_R(R^I)$.
    \end{enumerate}
\end{remark}
Next, we will talk about $I$-Ulrich modules introduced by Dao-Maitra-Sridhar \cite{Dao-Maitra-Sridhar}. The following definition is an equivalent condition for being an $I$-Ulrich module \cite[Theorem 4.6]{Dao-Maitra-Sridhar}. 
\begin{definition}
    Let $R$ be a one dimensional Cohen-Macaulay local ring and $I$ a regular ideal. Then a maximal Cohen-Macaulay $R$-module is called \textit{$I$-Ulrich} if there is an isomorphism $IM \cong M$. We denote by $\mathrm{Ul}_I(R)$ the category of $I$-Ulrich modules.
\end{definition}
\begin{remark}\label{i-ulrich-modules}
    Let $R$ be a one dimensional Cohen-Macaulay local ring and $I$ be a regular ideal. 
    \begin{enumerate}
        \item A module $M$ is $I$-Ulrich if and only if $M \in \CM(B(I))$ \cite[Theorem 4.6]{Dao-Maitra-Sridhar}. This has the consequence that if $M$ is an $I$-Ulrich module, then one has an inclusion $\tr(M) \subseteq b(I)$ and the converse holds when $M$ is reflexive.
        \item If $M$ is an $I$-Ulrich module, then for any $X \in \CM(R)$, the module $\Hom_R(M,X)$ is also an $I$-Ulrich module \cite[Lemma 4.15]{Dao-Maitra-Sridhar}.
        \item If $R$ is analytically unramified, then the normalisation $\overline{R}$ is $I$-Ulrich.
    \end{enumerate}
\end{remark}
Recall that a Cohen-Macaulay local ring $R$ is generically Gorenstein if for every minimal prime $p$ of $R$ the local ring $R_p$ is Gorenstein. In this case, the ring $R$ admits a canonical module $\omega$. In fact, the generically Gorenstein condition is equivalent to saying that the canonical module $\omega$ is isomorphic to an ideal. When we take $\omega$ to be an ideal we will refer to it as a canonical ideal. 
\begin{remark}\label{omega-Ulrich-remarks}
    Let $R$ be a one dimensional generically Gorenstein local ring with canonical ideal $\omega$.
    \begin{enumerate}
        \item There is an equality $\reflexive(R) = \Omega \CM(R)$. That is, an $R$-module $M$ is reflexive if and only if it is the syzygy of a maximal Cohen-Macaulay module.
        \item An $R$-module $M$ is $\omega$-Ulrich if and only if there is an isomorphism $M^* \cong DM$ \cite[Corollary 4.27]{Dao-Maitra-Sridhar}.
        \item For $n \gg 0$, the ideal $\omega^n$ is $\omega$-Ulrich. 
    \end{enumerate}
\end{remark}
The following theorem follows from \cite[Theorem 5.5]{Dao-Dey-Dutta}. We note that by the \textit{additive closure} $\add_R(X)$ of an $R$-module $X$, we mean the smallest subcategory that contains $X$ and is closed under direct summands, isomorphisms and finite direct sums.  
\begin{theorem}\label{main-theorem}
     Let $R$ be a one dimensional Cohen-Macaulay local ring which is analytically unramified. Let $\omega$ be a canonical ideal. Then, the following are equivalent.
    \begin{enumerate}
        \item The blowup algebra $B(\omega)$ is the normalisation $\overline{R}$.
        \item There is an equality $b(\omega) = \co(R)$.
        \item There exists an $n$ such that $\omega^n \cong \co(R)$.
    \end{enumerate}
    These equivalent conditions imply the following.
    \begin{enumerate}
        \item[(4)] There is an equality $\mathrm{Ul}_{\omega}(R) = \add_R(\overline{R})$. 
    \end{enumerate}
    When $R$ is a Henselian domain, all four conditions are equivalent.
\end{theorem}
We end this section with our definition and an example.
\begin{definition}
    Let $R$ be a one dimensional Cohen-Macaulay local ring with canonical ideal $\omega$. Assume that $R$ is analytically unramified and is not regular. We say that $R$ has \textit{minimal canonical conductor} if the equivalent conditions of Theorem \ref{main-theorem} hold for $R$. 
\end{definition}
\begin{example}\label{one-step-normal-example}
    Let $n \geq 3$ be a positive integer and $R = k[[t^n, t^{n+1}, \ldots, t^{2n-1}]]$. Then, by direct computation we can see that $\co(R)=\m$ and therefore one has equalities 
    \begin{align*}
        \co(R) = b(\omega) = \tr(\omega) = \m.
    \end{align*}
    Therefore, $R$ has minimal canonical conductor. Rings with maximal conductor have been considered by several authors in the last decade. For instance, they appear in Faber's work \cite{Faber} as \textit{one-step normal} rings. In the work of Dao-Dey-Dutta \cite{Dao-Dey-Dutta}, such rings are exactly one dimensional \textit{Ulrich split} rings. In the language of Herzog-Kumashiro-Stamate \cite{Herzog-Kumashiro-Stamate}, they are \textit{far-flung Gorenstein and nearly Gorenstein}. Recently, such rings also appeared in the work of Geroldinger-Yan-Zhang as \textit{conductor domains} \cite{Geroldinger-Yan-Zhong}.
\end{example}

    \section{Minimal canonical conductor versus far-flung Gorenstein}
    In Example \ref{one-step-normal-example}, we have seen that the canonical conductor $b(\omega)$ equals the canonical trace ideal $\tr(\omega)$. As we always have the inclusions $\co(R) \subseteq b(\omega) \subseteq \tr(\omega)$ for a one dimensional analytically unramified local ring, the condition $\tr(\omega) = \co(R)$ implies the condition $b(\omega) = \co(R)$. Hence, far-flung Gorenstein rings have minimal canonical conductor. However, not all rings with minimal canonical conductor are far-flung Gorenstein. We will see examples in the next section. in this section, we will look at the case when the two classes agree. We start with a definition due to Kumashiro \cite{Kumashiro}.
\begin{definition}
    Let $R$ be a generically Gorenstein Cohen-Macaulay local ring. The infimum of all nonnegative integers $n$ such that there exists a canonical ideal $\omega$ and an almost reduction $(a)$ of $\omega$ with $\omega^{n+1} = a\omega^n$ is called the \textit{canonical reduction number} of $R$ and is denoted by $\mathrm{can.red}(R)$.   
\end{definition}
As mentioned above, a one dimensional analytically unramified Cohen-Macaulay local ring is reduced and therefore it is generically Gorenstein.
\begin{proposition}
    Let $R$ be a one dimensional Cohen-Macaulay local ring which is analytically unramified. The following are equivalent.
    \begin{enumerate}
        \item $R$ has minimal canonical conductor and canonical reduction number 2.
        \item $R$ is far-flung Gorenstein.
    \end{enumerate} 
\end{proposition}
\begin{proof}
    The inclusion $b(\omega) \subseteq \tr(\omega)$ is an equality if and only if there exists an isomorphism $\tr(\omega) \cong \omega^*$ by \cite[Proposition 4.7]{Dao-Lindo}. And such an isomorphism holds if and only if $R$ has canonical reduction number at most 2 by \cite[Theorem 2.13]{Kumashiro}. This shows that if $R$ has minimal canonical conductor and canonical reduction number 2, then it is far-flung Gorenstein. On the other hand, as mentioned above all far-flung Gorenstein rings have minimal canonical conductor. Moreover, by \cite[Theorem 2.5]{Herzog-Kumashiro-Stamate}, we see that $R$ has canonical reduction number 2 when it is far-flung Gorenstein. This finishes the proof.
\end{proof}
We finish this section with an observation on the endomorphism ring $\m \colon \m$. This ring is often useful when one studies reflexive modules over $R$ and is far-flung Gorenstein whenever $R$ is. The latter assertion was proved by Herzog-Kumashiro-Stamate in \cite[Theorem 3.3]{Herzog-Kumashiro-Stamate} and the proof of the proposition below works similarly by making use of Claim 1 in the proof of \textit{op. cit}.
\begin{proposition}
    Let $(R, \m)$ be a one dimensional analytically unramified local Henselian domain. If $R$ has minimal canonical conductor, then the ring $B = \m \colon \m$ also has minimal canonical conductor. 
\end{proposition}


\section{Reflexive and Gorenstein birational extensions}
Our definition of a ring with minimal canonical conductor via the chain of ideals 
\begin{align*}
\co(R) \subseteq b(\omega) \subseteq \tr(\omega) \subseteq \m
\end{align*}
suggests that such a ring should be considered as far away as possible from being almost Gorenstein. However, this claim uses a specific characterisation of one dimensional almost Gorenstein rings, namely \cite[Proposition 7.1]{Dao-Maitra-Sridhar}. However, One of the most well-known characterisations of a one dimensional almost Gorenstein local ring $(R,\m)$ is via the endomorphism ring $E=\m \colon \m \cong \End_R(\m)$. Precisely, in \cite[Theorem 5.1]{Goto-Matsuoka-Phuong}, Goto-Matsuoka-Phuong proves that the following are equivalent for a one dimensional Cohen-Macaulay local ring $(R,\m)$:
\begin{enumerate}
    \item $E$ is a Gorenstein ring.
    \item $R$ is an almost Gorenstein ring with minimal multiplicity.
\end{enumerate}
In this section, we will explore this point of view. We start with some important remarks.
\begin{remark}
    Let $(R,\m)$ be a one dimensional Cohen-Macaulay local ring which is generically Gorenstein and put $E = \End_R(\m)$ as above. Then, both $\m$ and $E$ are reflexive $R$-modules, see for instance \cite[Lemma 2.5]{Kobayashi}. Assume that $R$ is not regular. Then, by \cite[Lemma 2.2]{Kobayashi-self-dual}, for any module-finite reflexive birational extension $S$ of $R$, one has a chain of inclusions 
    \begin{align*}
    R \subsetneq E \subseteq S.
    \end{align*}
    That is, $E$ is the minimal module-finite reflexive birational extension of $R$, in a sense. Combining this with the fact that any module-finite birational extension of a Gorenstein ring is also Gorenstein, we may rephrase the theorem of Goto-Matsuoka-Phuong as follows: for a one dimensional Cohen-Macaulay local ring $R$, the following are equivalent.
    \begin{enumerate}
        \item Any module-finite reflexive birational extension of $R$ is Gorenstein.
        \item $R$ is an almost Gorenstein ring with minimal multiplicity.
    \end{enumerate}      
\end{remark}
\begin{remark}
    Let $(R,\m)$ be a one dimensional Cohen-Macaulay local ring with infinite residue field and assume that it is not regular. In this case, $\m$ is not a principal ideal and therefore $\m$ is a regular proper trace ideal, see for instance \cite[Proposition 2.8]{Herzog-Rahimbeigi}. By \cite[Proposition 2.5]{Ooishi}, we have that $R$ has minimal multiplicity if and only if there exists an isomorphism $\m \cong \m^{-1} \cong \m^*$. And by \cite[Proposition 3.3]{Dao-Lindo}, the latter isomorphism is equivalent to saying that $\m$ is a stable trace ideal. Finally, let $\omega$ be the canonical module of $R$ and assume that $R$ is not Gorenstein so that $\omega$ is not free. If $R$ is almost Gorenstein, then by \cite[Theorem 7.1]{Dao-Maitra-Sridhar}, we have $b(\omega) = \m$. Combining all of these, we reach at the following paraphrase of the Goto-Matsuoka-Phuong theorem. The following are equivalent for a one dimensional Cohen-Macaulay local ring $R$ with canonical module $\omega$:
    \begin{enumerate}
        \item Any module-finite reflexive birational extension of $R$ is Gorenstein.
        \item $R$ is an almost Gorenstein ring and $b(\omega)$ is stable.
    \end{enumerate}
\end{remark}
Let $R$ be a one dimensional analytically unramified local ring. Then, the normalisation $\overline{R}$ of $R$ is a reflexive birational extension which is Gorenstein. Our main theorem in this section proves that this is the only such birational extension when $R$ has minimal canonical conductor.
\begin{theorem}\label{new-main-theorem}
     Let $R$ be a one dimensional Cohen-Macaulay local ring which is analytically unramified. Consider the two conditions.
    \begin{enumerate}
        \item $R$ has minimal canonical conductor.
        \item If $R \subseteq S \subseteq \overline{R}$ is a Gorenstein birational extension that is reflexive as an $R$-module, then $S = \overline{R}$. 
    \end{enumerate}
    Then, (1) implies (2). If $b(\omega)$ is stable, then (2) implies (1).
\end{theorem}
\begin{proof}
        Assume that $R$ has minimal canonical conductor. Let $S$ be a Gorenstein birational extension of $R$ which is reflexive as an $R$-module. Then, by \cite[Theorem 7.10]{Dao-Maitra-Sridhar}, we know that $\co(R,S)$ is $\omega$-Ulrich. Therefore, $\co(R,S)$ belongs to $\add_R(\co(R))$ and hence $\co(R,S) = \co(R)$ which shows that $S = \overline{R}$.

        On the other hand, assume (2) and that $b(\omega)$ is stable. We know that $B(\omega)$ is an $\omega$-Ulrich module. So, it is a reflexive birational extension of $R$. We know that $\tr(B(\omega))=\co(R, B(\omega)) = b(\omega)$ is $\omega$-Ulrich, as well. Moreover, by our assumption, it is also stable. Thus, by \cite[Theorem 7.10]{Dao-Maitra-Sridhar}, we see that $B(\omega)$ is Gorenstein. Hence, by our assumption, we get that $B(\omega) = \overline{R}$ and therefore, we are done by Theorem \ref{main-theorem}.
\end{proof}
We finish this section with an application to Arf local rings. The following definition and remarks on Arf local rings can be found in Lipman's original work on Arf rings \cite{Lipman}.
\begin{definition}
    Let $R$ be a one dimensional Noetherian semi-local Cohen-Macaulay ring.
    \begin{enumerate}
        \item We define a chain of blowup algebras inductively by letting $R_0 = R$ and $J_0 = J(R)$ (where $J(R)$ denotes the Jacobson radical of $R$) and setting $R_{n+1} := \End_{R_n}(J_n)$ where $J_n = J(R_n)$. If $\m$ is a maximal ideal of some $R_n$, then the local ring $(R_n)_\m$ is said to be \textit{infinitely near} to $R$.
        \item  We say that $R$ is an \textit{Arf ring} if every integrally closed ideal containing a nonzerodivisor has a principal reduction and whenever $x \in R$ is nonzerodivisor and $y,z \in R$ arbitrary elements with $y/x$ and $z/x$ in $\overline{R}$, we must have $yz/x \in R$.
    \end{enumerate}
\end{definition}
\begin{remark}
    The following are equivalent:
\begin{enumerate}
    \item $R$ is an Arf ring,
    \item Every integrally closed ideal containing a nonzerodivisor is stable,
    \item If $B$ is any local ring infinitely near to $R$, then $B$ has minimal multiplicity.
\end{enumerate}
It is clear from this that if $R = R_0$ is Arf, then $R_n$ is Arf for any $n \geq 0$. Moreover, if $R$ is a local ring with minimal multiplicity and $R_1$ is Arf, then $R$ is also Arf.
\end{remark}
Let $R$ be a complete local domain of dimension one and assume that it is analytically unramified. Then, we have a chain 
\begin{align*}
R = R_0 \subseteq R_1 \subseteq R_2 \subseteq \cdots \subseteq R_{n-1} \subseteq R_n = \overline{R}
\end{align*}
consisting of complete local domains which are infinitely near to $R$. By \cite[Theorem A]{Dao}, the $R_i$'s give a complete list of indecomposable reflexive $R$-modules up to isomorphism when $R$ is Arf. Now, assume that $R$ is Arf. If $R \subseteq S \subseteq \overline{R}$ is a birational extension that is reflexive as an $R$-module, then $S = R_i$ for some $i = 1, \ldots, n$. Hence, any birational extension $S \subseteq \overline{R}$ which is reflexive as an $R$-module is Arf. On the other hand, by \cite[Corollary 4.3]{Barucci-Froberg-value-semigroups}, $S$ is Gorenstein and Arf if and only if $S$ has multiplicity 2. This gives us the following characterisation by applying Theorem \ref{new-main-theorem}.
\begin{theorem}
    Let $R$ be a complete local Arf domain and assume that it is analytically unramified. The following are equivalent.
    \begin{enumerate}
        \item $R$ has minimal canonical conductor.
        \item For any birational extension $R\subseteq S \subseteq \overline{R}$ that is reflexive as an $R$-module, we have either $S = \overline{R}$ or $e(S) \geq 3$.
        \item There does not exist an $A_n$-singularity that is infinitely near to $R$.
    \end{enumerate}
\end{theorem}

\section{Numerical semigroups with minimal canonical conductor}
In this section, we classify numerical semigroup rings with minimal canonical conductor. By a numerical semigroup $H$, we mean a subsemigroup of $\mathbb{N}$ with $\mathbb{N}\setminus H$ finite. This means that there exists a largest number $f$ which does not belong to $H$. This $f$ is called the \textit{Frobenius} number of $H$. It has the property that $f + h \in H$ for any $h \in H$ with $h \neq 0$. This property can also hold for other numbers which do not belong to $H$. We put 
\begin{align*}
\mathrm{PF}(H):= \{x \in \mathbb{Z} \setminus H \colon x + h \in H \text{ for all } h \in H \text{ with } h \neq 0 \}
\end{align*}
and call an element of this set a \textit{pseudo-Frobenius} number. The Frobenius number $f$ is the largest pseudo-Frobenius number.

Given a field $k$, we denote by $k[[H]]$ the subring of the polynomial ring $k[[t]]$ generated by $t^h$ such that $h \in H$. There exists a nice interplay between $H$ and $k[[H]]$. For instance, every numerical semigroup has a minimal set of generators. The number of minimal generators of $H$ coincides with the embedding dimension of $k[[H]]$. The smallest element in $H$ coincides with the multiplicity of $k[[H]]$. Hence, $k[[H]]$ is of minimal multiplicity if and only if $H$ is minimally generated by elements of the form $n< a_2< \cdots < a_n$.

We can also see the type of $R = k[[H]]$ from the semigroup data. It is given by the size of $\mathrm{PF}(H)$. More precisely, if $f$ is the Frobenius number of $H$, then the $R$-module 
\begin{align*}
\omega = \sum_{x \in \mathrm{PF}(H)} Rt^{f-x}
\end{align*}
is a canonical module which is fractional ideal sitting between $R$ and $k[[t]] = \overline{R}$ \cite[Korollar 5]{Jager}.
 




\begin{theorem}\label{pseudo-frobenius-theorem}
    Let $H$ be a numerical semigroup with Frobenius number $f$. Then, the ring $R = k[[H]]$ has minimal canonical conductor if and only if $f-1$ is a pseudo-Frobenius number for $H$. 
\end{theorem}
\begin{proof}
    If $f-1$ is a pseudo-Frobenius number for $H$, then there exists a canonical fractional ideal $\omega = R + Rt + J$ where $J = R^{t_1} + \ldots + R^{t_s}$ for some $t_1, \ldots, t_s$. Note that $\omega^n \subseteq k[[t]] = \overline{R}$ for any positive integer $n$. On the other hand, we have 
    \begin{align*}
    (R+R^t)^n = R + Rt + Rt^2 + \cdots + Rt^n \subseteq \omega^n.
    \end{align*}
    For large enough $n$ (we can choose $n \leq f$), we have $(R+Rt)^n = k[[t]] = \overline{R}$. Hence, we obtain an isomorphism $\omega^n \cong \overline{R}$.

    If $f-1$ is not a pseudo-Frobenius number, we the canonical fractional ideal is of the form $\omega = R + R^{t_1} + \cdots + R^{t_s}$ for some $t_1, \ldots, t_s \geq 2$. Therefore, we have $t \notin \omega^n$ for any $n \geq 1$. This finishes the proof. 
\end{proof}
\begin{corollary}
    Let $H$ be a numerical semigroup minimally generated by $n< a_2< \cdots < a_n$. Then, the ring $k[[H]]$ has minimal canonical conductor if and only if $a_{n} = a_{n-1} +1$.
\end{corollary}
\begin{proof}
    Since the minimal number of generators equals the smallest generator, $k[[H]]$ is of minimal multiplicity. Therefore, pseudo-Frobenius numbers of $H$ are given by $a_2 - n, a_3-n, \ldots, a_n - n$. Hence, the Frobenius number is $a_n-n$ and $a_n-n-1$ is a Frobenius number if and only if $a_{n-1} = a_n - 1$.   
\end{proof}
\begin{example}
    Let $N$ be the semigroup generated by $n, n+1, n^2-n-1$ with $n \geq 3$ and $R = k[[N]]$ be the local semigroup ring generated by $N$. This ring was studied in \cite[Example 3.13]{Kumashiro}. It has a canonical fractional ideal $\omega = R + Rt$, it is nearly Gorenstein for all $n \geq 3$ and it is almost Gorenstein if and only if $n = 3$. We see from Theorem \ref{pseudo-frobenius-theorem} that 
    \begin{align*}
        \co(R) = b(\omega) \subseteq \tr(\omega) = \m.
    \end{align*}
    Therefore, $R$ has minimal canonical conductor. The inclusion $b(\omega) \subseteq \tr(\omega)$ is strict if and only if $n > 3$. Hence this ring, while being nearly Gorenstein, is far from being almost Gorenstein for $n > 3$. 
\end{example}

\iffalse
\section{Arf rings}
In this section, we focus on a class of rings which has been first studied by Lipman \cite{Lipman} and which has once again gained popularity in the last decade within commutative algebra. We start with some definitions.

\begin{definition}
    \begin{enumerate}
        \item Let $R$ be a commutative ring and $I$ an ideal which contains a nonzerodivisor. Then, $I$ is called \textit{stable} if $R^I = I \colon I\cong \End_R(I)$.
        \item Let $R$ be a one dimensional Noetherian semi-local Cohen-Macaulay ring. We define a chain of blowup algebras inductively by letting $R_0 = R$ and $J_0 = J(R)$ (where $J(R)$ denotes the Jacobson radical of $R$) and setting $R_{n+1} := \End_{R_n}(J_n)$ where $J_n = J(R_n)$. If $\m$ is a maximal ideal of some $R_n$, then the local ring $(R_n)_\m$ is said to be \textit{infinitely near} to $R$.
        \item  Let $R$ be a one dimensional Noetherian semi-local Cohen-Macaulay ring. We say that $R$ is an \textit{Arf ring} if every integrally closed ideal containing a nonzerodivisor has a principal reduction and whenever $x \in R$ is nonzerodivisor and $y,z \in R$ arbitrary elements with $y/x$ and $z/x$ in $\overline{R}$, we must have $yz/x \in R$.
    \end{enumerate}
\end{definition}
\begin{remark}
    The following are equivalent:
\begin{enumerate}
    \item $R$ is an Arf ring,
    \item Every integrally closed ideal containing a nonzerodivisor is stable,
    \item If $B$ is any local ring infinitely near to $R$, then $B$ has minimal multiplicity.
\end{enumerate}
It is clear from this that if $R = R_0$ is Arf, then $R_n$ is Arf for any $n \geq 0$. Moreover, if $R$ has is a local ring with minimal multiplicity and $R_1$ is Arf, then $R$ is also Arf.
\end{remark}
Recently, it was proved independently by Dao \cite{Dao}, Dao-Lindo \cite{Dao-Lindo} and Isobe-Kumashiro \cite{Isobe-Kumashiro} that under reasonable assumptions, an Arf ring has finite reflexive type: it admits finitely many indecomposable reflexive modules up to isomorphism. Below, we include a more precise version that appears in Dao's preprint and it also contains another result from that preprint which we will use later.
\begin{lemma}
    \begin{enumerate}
        \item Let $R$ be a semi-local, reduced Arf ring of dimension one which is complete with respect to its Jacobson radical. Then, the set of indecomposable reflexive $R$-modules is precisely the set of local rings near to $R$, up to isomorphism.
        \item Let $R$ be a semi-local, Cohen-Macaulay, equidimensional ring of dimension one. Assume that the completion of $R$ with respect its Jacobson radical is reduced. Then, $R$ is Arf if and only if any reflexive $R$-module $M$ is self-dual in the sense that $M \cong M^*$.
    \end{enumerate}
\end{lemma}

\section{ToDo}

\begin{enumerate}
    \item Make the definition for more general rings.
    \item Here is a new theorem which can be the new main theorem. 
    \begin{theorem}
    Let $R$ be a one dimensional Cohen-Macaulay local ring which is analytically unramified. Consider the two conditions.
    \begin{enumerate}
        \item $R$ has minimal canonical conductor.
        \item The only Gorenstein birational extension that is reflexive as an $R$-module is the normalisation $\overline{R}$ of $R$. 
    \end{enumerate}
    Then, (1) implies (2). If $b(\omega)$ is stable, then (2) implies (1).
    \end{theorem}
    \begin{proof}
        Assume that $R$ has minimal canonical conductor. Let $S$ be a Gorenstein birational extension of $R$ which is reflexive as an $R$-module. Then, by \cite[Theorem 7.10]{Dao-Maitra-Sridhar}, we know that $\co(R,S)$ is $\omega$-Ulrich. Therefore, $\co(R,S)$ belongs to $\add_R(\co(R))$ and hence $\co(R,S) = \co(R)$ which shows that $S = \overline{R}$.

        On the other hand, assume (b). We know that $B(\omega)$ is an $\omega$-Ulrich module. So, it is a reflexive birational extension of $R$. We know that $\tr(B(\omega))=\co(R, B(\omega)) = b(\omega)$ is $\omega$-Ulrich, as well. Moreover, by our assumption, it is also stable. Thus, by \cite[Theorem 7.10]{Dao-Maitra-Sridhar}, we see that $B(\omega)$ is Gorenstein. Hence, by our assumption, we get that $B(\omega) = \overline{R}$ and therefore, we are done.
    \end{proof}
\end{enumerate}
\fi

\bibliographystyle{alpha}
\bibliography{mcc.bib}
\end{document}
